\clearpage
\begin{proofbox}
\textbf{方法一：最低门槛——直接用两点斜率公式}
\begin{center}\begin{tikzpicture}[font=\small,>=stealth,line join=round]
\draw[gray!45,rounded corners=3pt] (0,0) rectangle (14.2,9.4);
\node at (7.1,8.8) {两条弦斜率相反，第三条弦的方向固定};
\begin{scope}[shift={(2.4,5.0)},scale=.85]
\draw[->,gray] (-.65,0)--(5.5,0) node[above] {$x$};
\draw[->,gray] (0,-4.8)--(0,3.3) node[left] {$y$};
\draw[teal!65!black,thick,samples=120,domain=-4.5:3.1,variable=\t] plot ({\t*\t/4},\t);
\draw[blue!75!black,very thick] (0,0)--(1,2);
\draw[orange!85!black,very thick] (1,2)--(4,-4);
\draw[green!45!black,very thick] (0,0)--(4,-4);
\fill (1,2) circle(2pt) node[above left=10pt] {$A$};
\fill (0,0) circle(2pt) node[above left=9pt] {$M=O$};
\fill (4,-4) circle(2pt) node[above right=12pt] {$N$};
\fill (1,0) circle(1.8pt) node[below right=10pt] {$F$};
\end{scope}
\node[anchor=west,align=left] at (8.7,5.3)
{$y^2=4x,\quad p=2$\\[12pt]$A(1,2),\quad F(1,0)$\\[12pt]
$M(0,0),\quad N(4,-4)$\\[16pt]
$\textcolor{blue!75!black}{k_{AM}=2},\quad\textcolor{orange!85!black}{k_{AN}=-2}$\\[16pt]
$\textcolor{green!45!black}{k_{MN}=-1}$};
\node[text width=13.2cm,align=center] at (7.1,.65) {坐标轴单位长度相同；斜率相反约束的是两条直线，并非两点坐标。};
\end{tikzpicture}
\end{center}


设
\[
 M=\left(\frac{u^2}{2p},u\right),\qquad
 N=\left(\frac{v^2}{2p},v\right),
\]
其中 $u\neq v$，$u\neq p$，$v\neq p$，且题设保证 $AM,AN$ 的斜率存在。

\textbf{第 1 步：求 $AM,AN$ 的斜率。} 因为 $A=(\frac p2,p)$，
\[
\begin{aligned}
 k_{AM}&=\frac{u-p}{\frac{u^2}{2p}-\frac p2}
 =\frac{2p(u-p)}{(u-p)(u+p)}=\frac{2p}{u+p},\\[3pt]
 k_{AN}&=\frac{2p}{v+p}.
\end{aligned}
\]
这里约去 $u-p$、$v-p$ 合法，因为 $M,N$ 均不等于 $A$；题设还保证 $u+p\ne0$、$v+p\ne0$。

\textbf{第 2 步：用斜率相反求纵坐标的和。}
\[
 \frac{2p}{u+p}=-\frac{2p}{v+p}
 \quad\Longrightarrow\quad \boxed{u+v=-2p}.
\]
\textbf{第 3 步：求弦 $MN$ 的斜率。} 由于 $u\ne v$，且 $u+v=-2p\ne0$，直线 $MN$ 的斜率存在，于是
\[
\begin{aligned}
 k_{MN}&=\frac{v-u}{\frac{v^2-u^2}{2p}}
 =\frac{2p}{u+v}
 =\boxed{-1}.
\end{aligned}
\]
\textbf{关键：}从两条斜率条件得到 $u+v$，再代入第三条弦；没有必要解出 $u,v$。
\end{proofbox}

\clearpage
\begin{proofbox}
\textbf{方法二：最快——斜率倒数恒等式（一行结束）}
\begin{center}\begin{tikzpicture}[font=\small,>=stealth,line join=round]
\draw[gray!45,rounded corners=3pt] (0,0) rectangle (14.2,9.4);
\node at (7.1,8.8) {同一个三角形：倒斜率给出横移与纵移之比};
\begin{scope}[shift={(2.4,5.0)},scale=.85]
\draw[->,gray] (-.65,0)--(5.5,0) node[above] {$x$};
\draw[->,gray] (0,-4.8)--(0,3.3) node[left] {$y$};
\draw[teal!65!black,thick,samples=120,domain=-4.5:3.1,variable=\t] plot ({\t*\t/4},\t);
\draw[blue!75!black,very thick] (0,0)--(1,2);
\draw[orange!85!black,very thick] (1,2)--(4,-4);
\draw[green!45!black,very thick] (0,0)--(4,-4);
\fill (1,2) circle(2pt) node[above left=10pt] {$A$};
\fill (0,0) circle(2pt) node[above left=9pt] {$M=O$};
\fill (4,-4) circle(2pt) node[above right=12pt] {$N$};
\fill (1,0) circle(1.8pt) node[below right=10pt] {$F$};
\end{scope}
\node[anchor=west,align=left] at (8.7,5.3)
{$\dfrac1{k_{AM}}=\dfrac12$\\[20pt]
$\dfrac1{k_{AN}}=-\dfrac12$\\[20pt]
$\dfrac1{k_{MN}}=-1$\\[20pt]
$\dfrac12-\dfrac12-1=-1$};
\node[text width=13.2cm,align=center] at (7.1,.65) {三条弦使用同一组端点；定点纵坐标 $p$ 被相加两次，需减去 $1$。};
\end{tikzpicture}
\end{center}


对于抛物线上纵坐标为 $s,t$ 的两个\textbf{不同点}，只要斜率有限、非零，就有
\[
 k=\frac{2p}{s+t}\qquad\Longleftrightarrow\qquad
 \frac1k=\frac{s+t}{2p}.
\]
因此
\[
 \frac1{k_{AM}}=\frac{p+u}{2p},\qquad
 \frac1{k_{AN}}=\frac{p+v}{2p}.
\]
先相加并利用题设：
\[
 \frac{u+v}{2p}=
 \frac1{k_{AM}}+\frac1{k_{AN}}-1=0-1=-1.
\]
于是 $u+v=-2p\ne0$，说明 $k_{MN}$ 确实存在；随后
\[
 \boxed{\frac1{k_{MN}}=
 \frac1{k_{AM}}+\frac1{k_{AN}}-1=-1}
 \quad\Longrightarrow\quad\boxed{k_{MN}=-1}.
\]
\textbf{记忆式：}
\[
 \boxed{\frac1{k_{MN}}=\frac1{k_{AM}}+\frac1{k_{AN}}-1}.
\]
该式以本题的 $A(\frac p2,p)$ 为定点，且只在三条相关斜率有限、非零时直接书写。它可视为抛物线两点弦斜率关系的一次直接应用。
\end{proofbox}

\clearpage
\begin{proofbox}
\textbf{方法三：常规联立——两条直线分别与抛物线求交}
\begin{center}\begin{tikzpicture}[font=\small,>=stealth,line join=round]
\draw[gray!45,rounded corners=3pt] (0,0) rectangle (14.2,7.2);
\node at (7.1,6.6) {分别找直线的第二交点，纵坐标中的变动部分相消};
\begin{scope}[shift={(1.4,3.8)},scale=.5]
\draw[->,gray] (-.5,0)--(9.7,0) node[above] {$x$};
\draw[->,gray] (0,-4.8)--(0,3.7) node[left] {$y$};
\draw[teal!65!black,thick,samples=110,domain=-4.3:3.4,variable=\t] plot ({\t*\t/2},\t);
\draw[gray,dashed] (-.2,1)--(4.5,1);
\draw[blue!75!black,very thick] (.5,1)--(2,2);
\draw[orange!85!black,very thick] (.5,1)--(8,-4);
\draw[green!45!black,very thick] (2,2)--(8,-4);
\fill (.5,1) circle(2pt) node[above left=10pt] {$A$};
\fill (2,2) circle(2pt) node[above=12pt] {$M$};
\fill (8,-4) circle(2pt) node[above right=12pt] {$N$};
\end{scope}
\node[anchor=west,align=left] at (8.8,3.8)
{$p=1,\quad t=\frac23$\\[14pt]
$u=\dfrac{2p}{t}-p=2$\\[16pt]
$v=-\dfrac{2p}{t}-p=-4$\\[16pt]
$u+v=-2p=-2$};
\node[text width=13.2cm,align=center] at (7.1,.65) {蓝、橙线分别与曲线在 $A$ 之外再交于 $M,N$；绿色弦斜率为 $-1$。};
\end{tikzpicture}
\end{center}


记 $k_{AM}=t$。抛物线上两不同点构成的非竖直弦斜率不为 $0$，故 $t\ne0$。题设给出 $k_{AN}=-t$。

\textbf{第 1 步：过 $A$ 的第一条直线。}
\[
 AM:\quad y-p=t\left(x-\frac p2\right).
\]
代入 $x=\frac{y^2}{2p}$，整理并因式分解：
\[
 2p(y-p)=t(y^2-p^2)
 \quad\Longleftrightarrow\quad
 (y-p)\,[\,2p-t(y+p)\,]=0.
\]
其中 $y=p$ 对应已知交点 $A$；另外一点 $M$ 的纵坐标为
\[
 \boxed{u=\frac{2p}{t}-p}.
\]
\textbf{第 2 步：第二条直线用斜率 $-t$。} 同理 $AN:y-p=-t(x-p/2)$，
代入抛物线得到另一交点 $N$ 的纵坐标
\[
 \boxed{v=-\frac{2p}{t}-p}.
\]
\textbf{第 3 步：相加，求第三条弦。}
\[
 u+v=\left(\frac{2p}{t}-p\right)
 +\left(-\frac{2p}{t}-p\right)=-2p,
\]
\[
 \boxed{k_{MN}=\frac{2p}{u+v}=-1}.
\]
\textbf{几何意义：}两条连线的斜率 $t,-t$ 可以变化，但产生的两个新交点纵坐标之和固定为 $-2p$，所以第三条弦始终\textbf{平行于直线 $y=-x$}。
\end{proofbox}


\clearpage
\begin{proofbox}
\begin{minipage}{\linewidth}
\textbf{补充证明：直接利用 C097，几步推出 C098}

记抛物线顶点为 $O$。先假设 $M,N$ 均不是 $O$，并且下面各条直线的斜率存在。
\begin{center}\begin{tikzpicture}[font=\small,>=stealth,line join=round]
\draw[gray!45,rounded corners=3pt] (0,0) rectangle (14.2,9.4);
\node at (7.1,8.8) {连接顶点 $O$，看清三次使用 C097 的连线};
\begin{scope}[shift={(2.7,4.8)}]
  \draw[->,gray] (-.8,0)--(3.5,0) node[above] {$x$};
  \draw[->,gray] (0,-3.65)--(0,3.1) node[left] {$y$};
  \draw[teal!65!black,thick,samples=120,domain=-3.45:2.8,variable=\t]
    plot ({\t*\t/4},\t);
  \coordinate (O) at (0,0);
  \coordinate (A) at (1,2);
  \coordinate (M) at (.25,-1);
  \coordinate (N) at (2.25,-3);
  \draw[gray!75,dashed,thick] (O)--(A) (O)--(M) (O)--(N);
  \draw[blue!75!black,very thick] (A)--(M);
  \draw[orange!85!black,very thick] (A)--(N);
  \draw[green!45!black,very thick] (M)--(N);
  \fill (O) circle(2pt) node[above left=7pt] {$O$};
  \fill (A) circle(2pt) node[above left=10pt] {$A$};
  \fill (M) circle(2pt) node[below=11pt] {$M$};
  \fill (N) circle(2pt) node[below left=10pt] {$N$};
\end{scope}
\node[anchor=west] at (7.1,7.7) {$y^2=4x,\quad p=2,\quad A(1,2)$};
\node[anchor=west] at (7.1,6.9) {$M(\tfrac14,-1),\quad N(\tfrac94,-3)$};
\node[anchor=west,text=blue!75!black] at (7.1,5.6)
  {$\dfrac1{k_{AM}}=\dfrac12-\dfrac14=\dfrac14$};
\node[anchor=west,text=orange!85!black] at (7.1,4.25)
  {$\dfrac1{k_{AN}}=\dfrac12-\dfrac34=-\dfrac14$};
\node[anchor=west,text=green!45!black] at (7.1,2.9)
  {$\dfrac1{k_{MN}}=-\dfrac14-\dfrac34=-1$};
\node[anchor=west] at (7.1,1.7) {$k_{OA}=2,\quad k_{OM}=-4,\quad k_{ON}=-\tfrac43$};
\node[text width=13.2cm,align=center] at (7.1,.65)
  {灰色虚线连接顶点 $O$；前两式相加时，$OA$ 的倒斜率被计入两次。};
\end{tikzpicture}
\end{center}
\end{minipage}

对图中三组连线应用 \textbf{C097（右开口抛物线两点弦的斜率倒数关系）}，得到
\[
\begin{aligned}
 \frac1{k_{AM}}&=\frac1{k_{OA}}+\frac1{k_{OM}},\\[3pt]
 \frac1{k_{AN}}&=\frac1{k_{OA}}+\frac1{k_{ON}},\\[3pt]
 \frac1{k_{MN}}&=\frac1{k_{OM}}+\frac1{k_{ON}}.
\end{aligned}
\]
\textbf{前两式相加，再利用第三式：}
\[
 \boxed{\frac1{k_{MN}}=\frac1{k_{AM}}+\frac1{k_{AN}}-\frac2{k_{OA}}}.
\]
由 $k_{AM}=-k_{AN}$，两项倒数相加为零，故
\[
 \frac1{k_{MN}}=-\frac2{k_{OA}}.
\]
原题 $A=(p/2,p)$，所以 $k_{OA}=\frac p{p/2}=2$，因此
\[
 \frac1{k_{MN}}=-1\quad\Longrightarrow\quad\boxed{k_{MN}=-1}.
\]
\textbf{顺手推广：}若把 $A$ 换成任意非顶点定点 $A_0=(\frac{a^2}{2p},a)$（$a\ne0$），上述三式仍成立。由于 $k_{OA_0}=\frac{2p}a$，所以
\[
 \boxed{\quad k_{MN}=-\frac12 k_{OA_0}=-\frac pa\quad}.
\]
\textbf{边界补充（不漏特殊点）：}C097 的上述倒数式不能直接用于 $M=O$ 或 $N=O$。例如 $M=O$ 时，$k_{A_0M}=\frac{2p}a$；令 $N=(\frac{v^2}{2p},v)$，由 $\frac{2p}{a+v}=-\frac{2p}a$ 得 $v=-2a$，从而 $k_{MN}=\frac{2p}{v}=-\frac pa$。$N=O$ 同理。另若 $MN$ 竖直则两点纵坐标之和为 $0$，而斜率对称必给出 $u+v=-2a\ne0$，矛盾。故本题的 $MN$ 斜率一定存在。
\end{proofbox}
